% Dijkstra 负边反例：s->a 2, s->b 5, b->a -4。算法定 a=2，真实距离 1（s->b->a）。
\begin{tikzpicture}[>=Stealth,
  v/.style={circle,draw=FigureStroke,fill=FigureBlueFill,line width=.8pt,minimum size=6.5mm,inner sep=0pt},
  vcur/.style={v,draw=FigureYellow,fill=FigureYellowFill},
  vopen/.style={v,fill=white},
  arc/.style={->,draw=FigureStroke,line width=.8pt},
  w/.style={fill=white,inner sep=.8pt,font=\figurefont\fontsize{7.5}{9}\selectfont},
  lab/.style={font=\figurefont\fontsize{7.5}{9}\selectfont},
  note/.style={font=\FigureNoteFont\fontsize{7.5}{9}\selectfont}]
% Panel labels share one baseline, 3 mm below the full row of content.
\begin{scope}[local bounding box=panel-content]
\begin{scope}
  \node[v] (s) at (0,0) {$s$};
  \node[vcur] (a) at (2,.9) {$a$};
  \node[vopen] (b) at (2,-.9) {$b$};
  \draw[arc,draw=FigureBlue,line width=1.4pt] (s)--node[w]{$2$}(a);
  \draw[arc] (s)--node[w]{$5$}(b);
  \draw[arc] (b)--node[w]{$-4$}(a);
  \node[lab,left=1mm of s] {$0$};
  \node[lab,right=1mm of a] {$d(a)=2$，确定};
  \node[lab,right=1mm of b] {$d(b)=5$};
  \coordinate[overlay] (panel-a-center) at (1.3,-1.85);
\end{scope}
\begin{scope}[xshift=5.4cm]
  \node[v] (s) at (0,0) {$s$};
  \node[v] (a) at (2,.9) {$a$};
  \node[vcur] (b) at (2,-.9) {$b$};
  \draw[arc] (s)--node[w]{$2$}(a);
  \draw[arc,draw=FigureGreen,line width=1.4pt] (s)--node[w]{$5$}(b);
  \draw[arc,draw=FigureGreen,line width=1.4pt] (b)--node[w]{$-4$}(a);
  \node[lab,left=1mm of s] {$0$};
  \node[lab,right=1mm of a,align=left] {仍为$2$\\{\color{FigureGreen}真实$\delta(a)=1$}};
  \node[lab,right=1mm of b] {$5$};
  \coordinate[overlay] (panel-b-center) at (1.3,-1.85);
\end{scope}
\end{scope}
\node[font=\figurefont\fontsize{8}{10}\selectfont,text=FigureInk,anchor=north,inner sep=0pt,outer sep=0pt] at ([yshift=-3mm]panel-a-center |- panel-content.south) {(a)};
\node[font=\figurefont\fontsize{8}{10}\selectfont,text=FigureInk,anchor=north,inner sep=0pt,outer sep=0pt] at ([yshift=-3mm]panel-b-center |- panel-content.south) {(b)};
\end{tikzpicture}
